\subsection{GNN}

\subsubsection{Anisotropic GNNs}
The difusco paper \cite{sun-23-difusco} relies on anisotropic GNNs. Those were initially introduced by \cite{bresson-18-nn-for-variable-graphs}. Instead of treating all neighbors equally, anisotropic GNNs allow to differentiate between neighbors - thus replacing summing, with different operators and using gates. LSTMs apparently proved to be too complex for the task. 
The proposed model is described as $$h_i^{\ell+1} = f_{\text{G-GCNN}}^{\ell} \left( h_i^{\ell} , \{h_j^{\ell} : j \to i\} \right) = \text{ReLU} \left( U^{\ell} h_i^{\ell} + \sum_{j \to i} \eta_{ij} \odot V^{\ell} h_j^{\ell} \right)$$
Meaning:
\begin{itemize}
    \item $h_{i}^{l+1}$: Node feature vector of node $i$ at layer $l+1$ (so defining the next step)
    \item $U^{l}h_{i}^{l}$: Multiply each node by the same matrix $U^{l}$ (learnable parameters)
    \item $\sum_{j\rightarrow i}$: Just aggregation over all neighbors of node $i$
    \item $V^{l}h_{j}^{l}$: Also a weight matrix, but it is multiplied with the neighbor features
    \item $\eta_{ij}$: The gate, which is computed as $\eta_{ij} = \sigma \left( A^{l} h_i^{l} + B^{l} h_j^{l} \right)$, with $A^{l}, B^{l}$ being learnable parameters again and $\sigma$ the sigmoid function. Note: \begin{itemize}
        \item A transforms the central node vector into a scalar value
        \item B does the same for the neighbor node. 
        \item Thus those two matrices decide how important the neighbor is for the central node.
    \end{itemize}
\end{itemize}

\question{Can other GNNs be used here? What do other papers say about the use of this GNN? What are the specific benefits of this architecture?}
